% GENERATED FILE. Edit canonical lessons, then run npm run generate:books.
% canonical-source-hash: fnv1a64:eefa46ceed89465c
% canonical-lessons: GU-C308-tarapo, GU-C308-sahas, GU-C308-vir, GU-C308-dhiraj, GU-C308-kargarvu, GU-R308-review-301-308

\chapter{A Brave Swimmer}
\label{ch:gu-a-brave-swimmer}
\hlchaptermodality{308}

\begin{chapteropening}
\textbf{By the end of this chapter:} \emph{I can tell how a brave swimmer saved a boy.}

Complete the last lesson of chapter 308: I can tell how a brave swimmer saved a boy.
\end{chapteropening}

\section[\texorpdfstring{\gu{તરાપો} (tarāpo)}{તરાપો (tarāpo)}]{\gu{તરાપો} (tarāpo) — a raft}

\label{lesson:GU-C308-tarapo}



\begin{quote}
Before the new one: say the Gujarati for safe and sound, then the Gujarati for phew.
\end{quote}

\subsection*{You'll want to know: \gu{તરાપો}}
\textbf{\gu{તરાપો}} — \emph{tarāpo} — \textquotedblleft{}a raft\textquotedblright{}.

\begin{quote}
\gu{પૂરમાં} \gu{ગામના} \gu{લોકોએ} \gu{એક} \gu{તરાપો} \gu{બનાવ્યો}.
\end{quote}

\begin{quote}
\emph{pūrmā̃ gāmnā lokoe ek tarāpo banāvyo.}
\end{quote}

\begin{quote}
During the flood the villagers built a raft.
\end{quote}

A masculine noun.

\begin{cousinweb}
\textbf{\gu{તરાપો}} (\emph{tarāpo}) comes from \emph{*tarpa}, \textquotedblleft{}a raft\textquotedblright{}, a form CDIAL reconstructs, and CDIAL lists the Gujarati noun, a raft of crossed timbers (CDIAL 5726). It is masculine: \textbf{\gu{તરાપો} \gu{બનાવ્યો}}, built a raft; \textbf{\gu{તરાપા} \gu{પર}}, on the raft.
\end{cousinweb}

\begin{grammarlens}[title={What you've built}]
One more word for telling what happened.
\end{grammarlens}

\subsection*{Your turn}
\begin{itemize}
\raggedright
  \item \emph{Say it:} \emph{tarāpo}
  \item \emph{Say it:} \emph{tarāpo}, once more
  \item \emph{Say it:} \emph{pūrmā̃ gāmnā lokoe ek tarāpo banāvyo.} — During the flood the villagers built a raft.
  \item \emph{Recall:} say the Gujarati for safe and sound, then the Gujarati for phew, then say \emph{tarāpo} again
\end{itemize}

\begin{tcolorbox}[breakable,colback=teal!4,colframe=teal!35!black,title={Before you move on}]
What does \gu{તરાપો} mean? (\textquotedblleft{}A raft\textquotedblright{}.) Say it once more.
\end{tcolorbox}
\section[\texorpdfstring{\gu{સાહસ} (sāhas)}{સાહસ (sāhas)}]{\gu{સાહસ} (sāhas) — daring; an adventure}

\label{lesson:GU-C308-sahas}



\begin{quote}
Before the new one: say the Gujarati for phew, then the Gujarati for a raft.
\end{quote}

\subsection*{You'll want to know: \gu{સાહસ}}
\textbf{\gu{સાહસ}} — \emph{sāhas} — \textquotedblleft{}daring; an adventure\textquotedblright{}.

\begin{quote}
\gu{નદીમાં} \gu{કૂદીને} \gu{અરુણે} \gu{મોટું} \gu{સાહસ} \gu{કર્યું}.
\end{quote}

\begin{quote}
\emph{nadīmā̃ kūdīne Aruṇe moṭũ sāhas karyũ.}
\end{quote}

\begin{quote}
By jumping into the river Arun showed great daring.
\end{quote}

A neuter noun.

\begin{cousinweb}
\textbf{\gu{સાહસ}} (\emph{sāhas}) is a tatsama, Sanskrit \emph{sāhasa}, \textquotedblleft{}boldness, rashness\textquotedblright{} (CDIAL 13378). Gujarati makes it good: daring, and the bold deed itself: \textbf{\gu{સાહસ} \gu{કર્યું}}.
\end{cousinweb}

\begin{grammarlens}[title={What you've built}]
One more word for telling what happened.
\end{grammarlens}

\subsection*{Your turn}
\begin{itemize}
\raggedright
  \item \emph{Say it:} \emph{sāhas}
  \item \emph{Say it:} \emph{sāhas}, once more
  \item \emph{Say it:} \emph{nadīmā̃ kūdīne Aruṇe moṭũ sāhas karyũ.} — By jumping into the river Arun showed great daring.
  \item \emph{Recall:} say the Gujarati for phew, then the Gujarati for a raft, then say \emph{sāhas} again
\end{itemize}

\begin{tcolorbox}[breakable,colback=teal!4,colframe=teal!35!black,title={Before you move on}]
What does \gu{સાહસ} mean? (\textquotedblleft{}Daring; an adventure\textquotedblright{}.) Say it once more.
\end{tcolorbox}
\section[\texorpdfstring{\gu{વીર} (vīr)}{વીર (vīr)}]{\gu{વીર} (vīr) — a hero; a brave man}

\label{lesson:GU-C308-vir}



\begin{quote}
Before the new one: say the Gujarati for a raft, then the Gujarati for daring.
\end{quote}

\subsection*{You'll want to know: \gu{વીર}}
\textbf{\gu{વીર}} — \emph{vīr} — \textquotedblleft{}a hero; a brave man\textquotedblright{}.

\begin{quote}
\gu{ગામના} \gu{લોકો} \gu{તેને} \gu{વીર} \gu{કહે} \gu{છે}.
\end{quote}

\begin{quote}
\emph{gāmnā loko tene vīr kahe che.}
\end{quote}

\begin{quote}
The villagers call him a hero.
\end{quote}

A masculine noun.

\begin{cousinweb}
\textbf{\gu{વીર}} (\emph{vīr}) is a tatsama, Sanskrit \emph{vīra}, \textquotedblleft{}a man, a hero\textquotedblright{} (CDIAL 12056). It is the brave one in war stories and in everyday courage alike.
\end{cousinweb}

\begin{grammarlens}[title={What you've built}]
One more word for telling what happened.
\end{grammarlens}

\subsection*{Your turn}
\begin{itemize}
\raggedright
  \item \emph{Say it:} \emph{vīr}
  \item \emph{Say it:} \emph{vīr}, once more
  \item \emph{Say it:} \emph{gāmnā loko tene vīr kahe che.} — The villagers call him a hero.
  \item \emph{Recall:} say the Gujarati for a raft, then the Gujarati for daring, then say \emph{vīr} again
\end{itemize}

\begin{tcolorbox}[breakable,colback=teal!4,colframe=teal!35!black,title={Before you move on}]
What does \gu{વીર} mean? (\textquotedblleft{}A hero; a brave man\textquotedblright{}.) Say it once more.
\end{tcolorbox}
\section[\texorpdfstring{\gu{ધીરજ} (dhīraj)}{ધીરજ (dhīraj)}]{\gu{ધીરજ} (dhīraj) — patience}

\label{lesson:GU-C308-dhiraj}



\begin{quote}
Before the new one: say the Gujarati for daring, then the Gujarati for a hero.
\end{quote}

\subsection*{You'll want to know: \gu{ધીરજ}}
\textbf{\gu{ધીરજ}} — \emph{dhīraj} — \textquotedblleft{}patience\textquotedblright{}.

\begin{quote}
\gu{મીરાએ} \gu{ધીરજથી} \gu{રાહ} \gu{જોઈ}.
\end{quote}

\begin{quote}
\emph{Mīrāe dhīrajthī rāh joī.}
\end{quote}

\begin{quote}
Mira waited patiently.
\end{quote}

A feminine noun.

\begin{cousinweb}
\textbf{\gu{ધીરજ}} (\emph{dhīraj}) probably goes back to Sanskrit \emph{dhairya}, \textquotedblleft{}firmness, calmness, patience\textquotedblright{} (Monier-Williams), the noun of \emph{dhīra}, \textquotedblleft{}steady\textquotedblright{}; CDIAL lists Gujarati \emph{dhīr}, \textquotedblleft{}patience\textquotedblright{}, under \emph{dhīra} but not this longer form (CDIAL 6817). It is kept: \textbf{\gu{ધીરજ} \gu{રાખો}}, be patient.
\end{cousinweb}

\begin{grammarlens}[title={What you've built}]
One more word for telling what happened.
\end{grammarlens}

\subsection*{Your turn}
\begin{itemize}
\raggedright
  \item \emph{Say it:} \emph{dhīraj}
  \item \emph{Say it:} \emph{dhīraj}, once more
  \item \emph{Say it:} \emph{Mīrāe dhīrajthī rāh joī.} — Mira waited patiently.
  \item \emph{Recall:} say the Gujarati for daring, then the Gujarati for a hero, then say \emph{dhīraj} again
\end{itemize}

\begin{tcolorbox}[breakable,colback=teal!4,colframe=teal!35!black,title={Before you move on}]
What does \gu{ધીરજ} mean? (\textquotedblleft{}Patience\textquotedblright{}.) Say it once more.
\end{tcolorbox}
\section[\texorpdfstring{\gu{કરગરવું} (kargarvũ)}{કરગરવું (kargarvũ)}]{\gu{કરગરવું} (kargarvũ) — to plead; to beg}

\label{lesson:GU-C308-kargarvu}



\begin{quote}
Before the new one: say the Gujarati for a hero, then the Gujarati for patience.
\end{quote}

\subsection*{You'll want to know: \gu{કરગરવું}}
\textbf{\gu{કરગરવું}} — \emph{kargarvũ} — \textquotedblleft{}to plead; to beg\textquotedblright{}.

\begin{quote}
\gu{છોકરાની} \gu{માતા} \gu{લોકોને} \gu{મદદ} \gu{માટે} \gu{કરગરતી} \gu{હતી}.
\end{quote}

\begin{quote}
\emph{chokrānī mātā lokone madad māṭe kargartī hatī.}
\end{quote}

\begin{quote}
The boy's mother was pleading with people for help.
\end{quote}

Past (she …): \gu{તે} \gu{માતા} \gu{કરગરી}, \emph{te mātā karagrī}.

\begin{cousinweb}
\textbf{\gu{કરગરવું}} (\emph{kargarvũ}) No dictionary we use records an older form, so its history is not settled. The person you plead with takes \textbf{\gu{ને}}: \textbf{\gu{લોકોને} \gu{કરગરતી} \gu{હતી}}.
\end{cousinweb}

\begin{grammarlens}[title={What you've built}]
One more word for telling what happened.
\end{grammarlens}

\subsection*{Your turn}
\begin{itemize}
\raggedright
  \item \emph{Say it:} \emph{kargarvũ}
  \item \emph{Say it:} \emph{kargarvũ}, once more
  \item \emph{Say it:} \emph{chokrānī mātā lokone madad māṭe kargartī hatī.} — The boy's mother was pleading with people for help.
  \item \emph{Recall:} say the Gujarati for a hero, then the Gujarati for patience, then say \emph{kargarvũ} again
\end{itemize}

\begin{tcolorbox}[breakable,colback=teal!4,colframe=teal!35!black,title={Before you move on}]
What does \gu{કરગરવું} mean? (\textquotedblleft{}To plead; to beg\textquotedblright{}.) Say it once more.
\end{tcolorbox}
\section[(dialogue)]{Looking back: from A Building on Fire to A Brave Swimmer}

\label{lesson:GU-R308-review-301-308}



\begin{quote}
Say the Gujarati for a large building and a ladder.
\end{quote}

\subsection*{Your turn}
Say these aloud:

\begin{itemize}
\raggedright
  \item a large building, a ladder, to choke, to burst, a bang
  \item a fire engine, to extinguish, a rope, to melt, a blaze
  \item a cyclone, to strike, a thatched roof, a hut, to lay waste
  \item to be swept away, to submerge, damage, a calamity, a crisis
  \item to shake, a heap, to give way, to be buried, rubble
  \item a volunteer, relief, a camp, grain, to throng
  \item to rescue, indebted, deep, safe and sound, phew
  \item a raft, daring, a hero, patience, to plead
\end{itemize}

\begin{tcolorbox}[breakable,colback=teal!4,colframe=teal!35!black,title={Before you move on}]
A large building? (\textbf{\gu{ઇમારત}}, \emph{imārat}.) A ladder? (\textbf{\gu{સીડી}}, \emph{sīḍī}.) To choke? (\textbf{\gu{ગૂંગળાવું}}, \emph{gū̃gḷāvũ}.) To burst? (\textbf{\gu{ફાટવું}}, \emph{phāṭvũ}.) A bang? (\textbf{\gu{ધડાકો}}, \emph{dhaḍāko}.) A fire engine? (\textbf{\gu{બંબો}}, \emph{bambo}.) To extinguish? (\textbf{\gu{બુઝાવવું}}, \emph{bujhāvvũ}.) A rope? (\textbf{\gu{દોરડું}}, \emph{dorḍũ}.) To melt? (\textbf{\gu{ઓગળવું}}, \emph{ogaḷvũ}.) A blaze? (\textbf{\gu{ભડકો}}, \emph{bhaḍko}.) A cyclone? (\textbf{\gu{વાવાઝોડું}}, \emph{vāvājhoḍũ}.) To strike? (\textbf{\gu{ત્રાટકવું}}, \emph{trāṭakvũ}.) A thatched roof? (\textbf{\gu{છાપરું}}, \emph{chāprũ}.) A hut? (\textbf{\gu{ઝૂંપડું}}, \emph{jhū̃pḍũ}.) To lay waste? (\textbf{\gu{ઉજાડવું}}, \emph{ujāḍvũ}.) To be swept away? (\textbf{\gu{તણાવું}}, \emph{taṇāvũ}.) To submerge? (\textbf{\gu{ડુબાડવું}}, \emph{ḍubāḍvũ}.) Damage? (\textbf{\gu{નુકસાન}}, \emph{nuksān}.) A calamity? (\textbf{\gu{આફત}}, \emph{āphat}.) A crisis? (\textbf{\gu{સંકટ}}, \emph{sankaṭ}.) To shake? (\textbf{\gu{હલવું}}, \emph{halvũ}.) A heap? (\textbf{\gu{ઢગલો}}, \emph{ḍhaglo}.) To give way? (\textbf{\gu{ધસવું}}, \emph{dhasvũ}.) To be buried? (\textbf{\gu{દટાવું}}, \emph{daṭāvũ}.) Rubble? (\textbf{\gu{કાટમાળ}}, \emph{kāṭmāḷ}.) A volunteer? (\textbf{\gu{સ્વયંસેવક}}, \emph{svayansevak}.) Relief? (\textbf{\gu{રાહત}}, \emph{rāhat}.) A camp? (\textbf{\gu{છાવણી}}, \emph{chāvṇī}.) Grain? (\textbf{\gu{અનાજ}}, \emph{anāj}.) To throng? (\textbf{\gu{ઊમટવું}}, \emph{ūmaṭvũ}.) To rescue? (\textbf{\gu{ઉગારવું}}, \emph{ugārvũ}.) Indebted? (\textbf{\gu{ઋણી}}, \emph{ṛṇī}.) Deep? (\textbf{\gu{ઊંડું}}, \emph{ū̃ḍũ}.) Safe and sound? (\textbf{\gu{હેમખેમ}}, \emph{hemkhem}.) Phew? (\textbf{\gu{હાશ}}, \emph{hāś}.) A raft? (\textbf{\gu{તરાપો}}, \emph{tarāpo}.) Daring? (\textbf{\gu{સાહસ}}, \emph{sāhas}.) A hero? (\textbf{\gu{વીર}}, \emph{vīr}.) Patience? (\textbf{\gu{ધીરજ}}, \emph{dhīraj}.) To plead? (\textbf{\gu{કરગરવું}}, \emph{kargarvũ}.)
\end{tcolorbox}
